
\documentclass[a4paper,11pt]{article}
\usepackage{fontspec}
\setmainfont{Noto Sans}
\usepackage[french]{babel}
\usepackage{microtype}
\usepackage[top=9mm,bottom=8mm,left=11mm,right=11mm]{geometry}
\usepackage{xcolor}
\usepackage{tikz}
\usetikzlibrary{arrows.meta,calc,positioning}
\usepackage[most]{tcolorbox}
\usepackage{ulem}
\pagestyle{empty}

\definecolor{Blue}{HTML}{356FA8}
\definecolor{BlueDark}{HTML}{244D78}
\definecolor{BlueLight}{HTML}{EAF2F9}
\definecolor{BlueBorder}{HTML}{B9D3E3}
\definecolor{Violet}{HTML}{7A5AA6}
\definecolor{VioletDark}{HTML}{563A7A}
\definecolor{VioletLight}{HTML}{F2ECF8}
\definecolor{VioletBorder}{HTML}{D3C4E2}
\definecolor{Neutral}{HTML}{6C7378}
\definecolor{NeutralDark}{HTML}{42494D}
\definecolor{NeutralLight}{HTML}{F0F2F3}
\definecolor{NeutralBorder}{HTML}{C9CED1}
\definecolor{Ink}{HTML}{172A32}
\definecolor{Grey}{HTML}{64757D}
\definecolor{Red}{HTML}{D52F35}
\definecolor{RedLight}{HTML}{FCEBED}

\tcbset{
 enhanced,boxrule=.6pt,arc=2.8mm,
 left=3.7mm,right=3.7mm,top=1.75mm,bottom=1.75mm,
 boxsep=1mm,coltext=Ink
}
\newtcolorbox{bluebox}[1][]{colback=BlueLight,colframe=BlueBorder,#1}
\newtcolorbox{violetbox}[1][]{colback=VioletLight,colframe=VioletBorder,#1}
\newtcolorbox{neutralbox}[1][]{colback=NeutralLight,colframe=NeutralBorder,#1}
\newtcolorbox{warnbox}[1][]{colback=RedLight,colframe=Red,
 left=2.7mm,right=2.7mm,top=1.25mm,bottom=1.25mm,#1}

\begin{document}
\begin{tikzpicture}[remember picture,overlay]
\draw[draw=NeutralBorder,line width=.9pt,rounded corners=1.7mm]
([xshift=1mm,yshift=-1mm]current page.north west)--([xshift=-1mm,yshift=-1mm]current page.north east)--
([xshift=-1mm,yshift=1mm]current page.south east)--([xshift=1mm,yshift=1mm]current page.south west)--cycle;
\fill[Neutral]([xshift=1mm,yshift=-1mm]current page.north west) rectangle ([xshift=5.2mm,yshift=1mm]current page.south west);
\end{tikzpicture}
\vspace*{1mm}\hspace*{6mm}{\color{NeutralDark}\fontsize{25}{27}\selectfont\bfseries PARTICIPI PASSAT}
\vspace{.5mm}\hspace*{6mm}{\color{Grey}\large Le participe passé}
\par\vspace{2.5mm}\noindent\begin{minipage}{.91\textwidth}
\begin{neutralbox}\textbf{\color{NeutralDark}FORMAS}\\[1mm]\centering
{\Large\textbf{PARTICIPI PASSAT DEL VÈRB}}\\[1mm]
\small Exemples : \textit{cantar $\rightarrow$ cantat \quad faire $\rightarrow$ fach \quad venir $\rightarrow$ vengut}
\end{neutralbox}
\vspace{2.5mm}\begin{center}\begin{tikzpicture}[>=Latex]
\draw[Ink!65,line width=.8pt](0,0)--(15.1,0);
\node[below=2mm]at(2.2,0){\small PASSAT};\node[below=2mm]at(7.5,0){\small ARA};\node[below=2mm]at(13.5,0){\small FUTUR};
\draw[Neutral,line width=2.5pt](2.5,0)--(5.4,0);
\draw[NeutralDark,line width=1.1pt,->](7.5,1)--(7.5,.18);
\node[above=3.8mm]at(7.5,1){\bfseries\color{NeutralDark}ARA};
\node[above=3mm,align=center,text width=6.2cm]at(4,0){\small\bfseries accion acabada / resultat};
\end{tikzpicture}\end{center}
\vspace{1.5mm}\begin{neutralbox}\textbf{\color{NeutralDark}A QUÉ SERVÍS ?}\\[1mm]
\begin{itemize}\setlength{\itemsep}{.4mm}
\item formar de nombrosas formas verbalas compostas ;
\item exprimir lo resultat o l'estat resultant d'una accion ;
\item èsser emplegat coma adjectiu dins certanas construccions ;
\item s'acordar en genre e en nombre quand la construccion o l'usatge o demanda.
\end{itemize}\end{neutralbox}
\vspace{2mm}\begin{neutralbox}\textbf{\color{NeutralDark}EXEMPLES}\\[1mm]
\textbf{Ai cantat tota la serada.}\\
\textbf{Es vengut fòrça lèu.}\\
\textbf{La pòrta es tampada.}\\
\textbf{Las letras son escritas.}
\end{neutralbox}
\vspace{2mm}
\begin{minipage}{.485\textwidth}\begin{neutralbox}[height=3.15cm]
\textbf{\color{NeutralDark}FORMAS IRREGULARAS}\\[1mm]
\centering
\textbf{èsser $\rightarrow$ estat}\\
\textbf{aver $\rightarrow$ agut}\\
\textbf{faire $\rightarrow$ fach}\\
\textbf{dire $\rightarrow$ dich}\\
\textbf{venir $\rightarrow$ vengut}
\end{neutralbox}\end{minipage}\hfill
\begin{minipage}{.485\textwidth}\begin{warnbox}[height=3.15cm]
\textbf{\color{Red}ATENCION}\\[1mm]
Lo participi passat pòt variar en \textbf{genre e nombre}.\\[1mm]
Ex. \textit{vengut / venguda / venguts / vengudas}.
\end{warnbox}\end{minipage}
\vspace{2mm}\begin{neutralbox}\textbf{\color{NeutralDark}A RETÉNER}\\[1mm]
Lo participi passat es una forma clau de la conjugason occitana : participa a las \textbf{formas compostas}
e pòt tanben aver una valor adjectivala.
\end{neutralbox}
\vspace{1.5mm}\centering{\small\color{Grey}\textbf{Exemples :} cantat, finit, vendut, fach, dich, escrit, vengut...}
\vspace{1mm}{\scriptsize\color{Grey}Referéncia : occitan languedocian, graphia classica, vèrb'Òc / Lo Congrès.}
\end{minipage}\end{document}
